\documentclass[12pt]{article}

\usepackage[utf8]{inputenc}
\usepackage{amsmath, amssymb, amsthm}
\usepackage{geometry}
\geometry{a4paper, margin=2.5cm}

\title{Algebră de Mulțimi}
\author{}
\date{}

\theoremstyle{definition}
\newtheorem{definition}{Definiție}

\begin{document}

\maketitle

\begin{definition}
Fie $\Omega$ un spațiu. O familie de submulțimi $\mathcal{A} \subseteq \mathcal{P}(\Omega)$ se numește \emph{algebră de mulțimi} dacă:
\begin{enumerate}
    \item $\Omega \in \mathcal{A}$.
    \item Dacă $A \in \mathcal{A}$, atunci $A^{c} \in \mathcal{A}$.
    \item Dacă $A, B \in \mathcal{A}$, atunci $A \cup B \in \mathcal{A}$.
\end{enumerate}
\end{definition}

\noindent
O algebră este închisă la complement și reuniuni finite.

\end{document}
